\documentclass{article}
\usepackage{amsmath,amssymb,amsthm}
\newtheorem{theorem}{Theorem}
\newtheorem{lemma}{Lemma}
\begin{document}
\title{A008: the original conditioned block-forest reference for unequal read-once DNFs}
\author{A008, B02}\maketitle
All costs and entropies are in bits. Natural logarithms occur only in the source rank statistic $F$. This is a reference-cost closure on a known FEI class, not a new Boolean class theorem.

Let $m\ge1$, let $R_1,\ldots,R_m$ be disjoint nonempty input blocks with widths $w_l\ge1$, and let $h_l$ be the sign of a conjunction of one signed literal at each coordinate in $R_l$. Let $f=h_1\vee\cdots\vee h_m$, true equal to $+1$. Constant clauses are removed as in the source. Input sign flips do not change squared Fourier coefficients or physical influences, so the proof takes all literals positive. Variables outside the blocks are irrelevant and omitted.

Set $p_l=2^{-w_l}$, $q_l=1-p_l$, $Q_0=\prod_lq_l$, $\pi_l=p_l$, $M_l=1$, and $\lambda_l=1$. Fix any one full order in each block. This is one globally fixed source queue: independently choose roots with odds $p_l$, choose subset size with probability $1/[s(s+1)]$, choose the rank-weighted subset, and choose child count with probability $1/[(c+1)(c+2)]$. Child-label choices use the source $\lambda$ rule. Every coordinate is in one block, so the fractional-incidence constraints hold. Put
\[
 \gamma=\frac65,\quad \Lambda=\sum_l\log(1+p_l),\quad
 Z=1-2^{-\Lambda},\quad D=2^\Lambda-1=\prod_l(1+p_l)-1.
\]

\begin{theorem}
For every such read-once DNF, the posterior-forest residual of the specified original prime-implicant table satisfies
\[
 \mathcal R_f\le\left(\frac{16}{5}+\log e\right)I(f)
 <\frac{47}{10}I(f).
\]
There is exactly one partition-forest run for every nonempty output set $A$. Its actual conditioned probability $q_{\rm pf}(A)$, including every subset/count normalizer, has spectral expected cost
\[
 C_{\rm pf}:=\sum_{A\ne\varnothing}P_f(A)\log\frac1{q_{\rm pf}(A)}
 \le\left(\frac{26}{5}+2\log3+2\log e\right)I(f).
\]
Hence the complete entropy-reference cost satisfies
\[
 h_2(\operatorname{Var}f)+C_{\rm pf}<\frac{67}{5}I(f).
\]
This last bound includes the empty/nonempty split. It does not include an optional extra binary-prefix rounding bit or an optional mixture selector. No selector is needed for the single fixed reference asserted here. It is stronger than applying the source's coarser generic $K_{\rm ref}$ to the residual.
\end{theorem}

\begin{proof}
Write $a_l=\mathbf1_{\{h_l=1\}}$, $b_l=1-a_l$ and $f=1-2\prod_lb_l$. For a nonempty set $A$, write $A_l=A\cap R_l$ and $K=\{l:A_l\ne\varnothing\}$. Each $b_l$ has coefficient $q_l$ at the empty set and coefficient of magnitude $p_l$ at every nonempty local subset. Thus
\[
 P_f(A)=4\prod_{l\in K}p_l^2\prod_{l\notin K}q_l^2,
 \qquad P_f(\varnothing)=(1-2Q_0)^2,
 \qquad v=4Q_0(1-Q_0).
\]
There are $2^{w_l}-1=q_l/p_l$ nonempty local subsets. Summing over all other blocks proves the exact hit mass
\[
 \mu_l:=P_f(A_l\ne\varnothing)=4p_lQ_0.
\]
Conditional on hitting block $l$, $A_l$ is uniform on its nonempty subsets. More sharply, conditional on a particular coordinate $j\in R_l$ being selected, every other coordinate of that block is an independent fair selection. Therefore
\[
 \mathbb E[|A_l|\mid A_l\ne\varnothing]=\frac{w_l}{2q_l},
 \qquad I_l=\mu_l\frac{w_l}{2q_l}
              =\frac{2Q_0w_lp_l}{q_l},\quad
 I=2Q_0\sum_l\frac{w_lp_l}{q_l}.
\]
Independently, a coordinate of block $l$ is physically pivotal when all its $w_l-1$ companion literals are true and all other clauses are false. This gives influence $2p_lQ_0/q_l$ per coordinate and confirms the same formula.

The clauses are precisely the prime implicants: an implicant containing no entire clause leaves at least one unfixed literal in each block; set one such literal false in each block to produce a false completion. Minimal implicants therefore are exactly the full individual clauses. Every selected coordinate has a unique gate owner. The used incidence graph has no edges, every used gate is an isolated source SCC and must be a root, and there are no children or attachments. The root list is in the fixed source order, with no permutation multiplicity. Subsets are $A_l$, and $s_l=\kappa_l=|A_l|$; hence the cover and congestion terms are exactly zero. There is one actual partition run per $A$, so posterior entropy is zero.

The probability of its initial roots, after conditioning on at least one root, is
\[
 \frac{\prod_{l\in K}p_l}{D}.
\]
Thus $\Lambda+\log Z=\log D$ is retained exactly, including all absent-root factors. The residual is
\[
 \mathcal R_f=v\log D+\sum_l\mu_lw_l
 +\gamma\mathbb E_{P_f}\sum_{j\in A}
       \log\frac{r_l(j)}{|A\cap U_j^l|+1}.
 \tag{1}
\]
Gate labels have $L_l=\log(1/\lambda_l)=0$, but gate selection is not free: it is the paid initial-root experiment, represented by $v\log D+\sum_l\mu_lw_l$. Each used gate name is specified once in that experiment.

For rarity,
\[
 \sum_l\mu_lw_l=4Q_0\sum_lw_lp_l\le2I,
\]
because $q_l\le1$. For normalization, write $S=\sum_lp_l$, $T=\sum_lw_lp_l/q_l$. Since $w_l/q_l\ge2$ for every positive integer width, $T\ge2S$. Also $\log D<\Lambda\le S\log e$. Consequently
\[
 v\log D\le4Q_0(1-Q_0)S\log e
 \le (\log e)I.
 \tag{2}
\]
This also follows from the source weighted-activation lemma; the argument above is a direct block proof. Negative $\log D$ cases remain valid and have not been replaced by absolute values.

For rank, let a selected coordinate have physical rank $r$ in its block. Its selected predecessor count is $X\sim\operatorname{Bin}(r-1,1/2)$. The binomial identity, proved by $\binom{r-1}{k}/(k+1)=\binom r{k+1}/r$, is
\[
 \mathbb E\frac{r}{1+X}=2(1-2^{-r})\le2.
\]
Jensen's inequality for the concave logarithm gives
\[
 \mathbb E\left[\left.\log\frac r{1+X}\right|j\in A\right]\le1.
\]
Multiply by $P_f(j\in A)$ and sum to obtain
\[
 \mathbb E\sum_{j\in A}\log\frac{r_l(j)}{|A\cap U_j^l|+1}\le I.
 \tag{3}
\]
Equations (1)--(3) prove $\mathcal R_f\le(2+\log e+6/5)I$. Since $\ln2>2/3$ by the positive atanh series, $\log e<3/2$, giving the strict $47I/10$ bound. The earlier factorial/concavity estimate is unnecessary in the final proof.

To audit the \emph{actual} complete probabilities, put
\[
 \theta_s^l(j)=\min\{1,(s/r_l(j))^\gamma\},
 \quad Z_s^l=e_s((\theta_s^l(j))_{j\in R_l}),
 \quad t_l(B)=\frac1{s(s+1)}\frac{\prod_{j\in B}\theta_s^l(j)}{Z_s^l}.
\]
The selected child count zero has probability $1/2$ at every used gate. Its attachment multiset is empty and has probability one. Therefore
\[
 q_{\rm pf}(A)=\frac1D\prod_{l\in K}\frac{p_l}{2}t_l(A_l).
 \tag{4}
\]
Each expression is one specified finite successful event of the original generator. These events are disjoint, so $\sum_{A\ne\varnothing}q_{\rm pf}(A)\le1$. Other successful runs, including repeated self-visits, may increase the original output mass; none is used here. Formula (4) contains the gate identities, all size probabilities, subset normalizers/weights and zero-child count probabilities. It specifies a terminating decoder: read the root set once, read one local subset per root in fixed order, and output their union. Disjoint supports identify the dictionary uniquely from the output, so there is no latent-witness ambiguity credit.

Its ideal expected cost is
\[
 C_{\rm pf}=v\log D+\sum_l\mu_l\left[w_l+1+
   \mathbb E_{B_l\text{ uniform, nonempty}}
   \left(\log(s_l(s_l+1))+\log Z_{s_l}^l
       -\sum_{j\in B_l}\log\theta_{s_l}^l(j)\right)\right].
 \tag{5}
\]
The root terms in (5) are at most $(2+\log e)I$. Child count zero costs $\sum_l\mu_l\le I$. Since $s(s+1)\le3^s$, size costs at most $(\log3)I$. For one order, $\sum_j\theta_s(j)\le6s$ by the integral tail of $r^{-6/5}$, and $e_s(\theta)\le(\sum\theta)^s/s!\le(6e)^s$. Thus subset normalizers cost at most $(\log(6e))I$. Finally, pointwise $|A\cap U_j^l|+1\le s_l$ and the rank ratio is at least one, so
\[
 -\log\theta_{s_l}^l(j)=\gamma\log^+\frac{r_l(j)}{s_l}
 \le\gamma\log\frac{r_l(j)}{|A\cap U_j^l|+1}.
\]
Equation (3) pays this by $\gamma I$. Summing gives
\[
 C_{\rm pf}\le[3+\gamma+\log e+\log3+\log(6e)]I
 =\left[\frac{26}{5}+2\log3+2\log e\right]I.
\]
The standard source inequality $h_2(v)\le2I$ and Gibbs applied to the subprobability reference with mass $1-v$ at empty and $v q_{\rm pf}(A)$ elsewhere give the claimed full bound. Since $3^5<2^8$, $\log3<8/5$; together with $\log e<3/2$, the full coefficient is strictly below $36/5+16/5+3=67/5$.
\end{proof}

\paragraph{Choice of reference and units.}
No spectrum-dependent table, gate subset law or scale is selected for free. The clause table, $\lambda_l=1$ and the within-block orders are fixed from the function's read-once DNF. Every root rarity is $w_l$ bits, and the normalizer uses the same vector. The source natural-rank factor $\gamma/\ln2$ becomes $\gamma$ times the base-two rank sum used above. Complementing the output gives the same spectrum and hence the same numerical statement for read-once CNFs through the complementary union representation.

\paragraph{Universal boundary.}
The proof does not say arbitrary spectra have uniform local subset laws. Its two decisive properties are the forced once-per-block assignment and the exact conditional binomial predecessor law. For overlapping prime implicants both may fail. The first unsupported PI-J line remains payment of the full joint label/rank/subset cost for an actual posterior under one globally fixed fractional vector. Root cancellation alone does not supply this property, and a coordinatewise affine-plane pressure obstruction is not a refutation of the whole-set bound.
\end{document}
